\scanpage{25}
\margin{24}
\begin{center}- 2 -\end{center}
\begin{flushright}
MR LE PROFESSEUR GROTHENDIECK ALEX.\\
MATHEMATIQUES \hfill 5
\end{flushright}

Ces documents doivent \uline{impérativement} me parvenir, datés et signés par chaque chef de projet, \uline{avant le 15 mai prochain} à l'adresse suivante :

\begin{center}
Ministère des Universités\\
Délégation aux Relations Universitaires\\
Internationales\\
Bureau des Echanges Scientifiques\\
173, boulevard Saint-Germain\\
75006 PARIS
\end{center}

Il peut s'agir :

1) De projets nouveaux pour lesquels les moyens sont accordés annuellement (les demandes ponctuelles ne relevant pas des projets intégrés).

2) De projets en cours. Il est rappelé que la demande de moyens doit être renouvelée chaque année. Outre le dossier ci-joint, vous voudrez bien me faire parvenir un bilan détaillé faisant apparaître l'état d'avancement du projet, les résultats obtenus et les perspectives de développement de la coopération..

Les responsables français sont tenus d'une part de présenter un projet, même si les responsables québécois, ont, de leur côté, établi un projet parallèle, d'autre part, de mentionner les noms et qualités des personnes impliquées.

J'appelle votre attention sur la possibilité d'accorder à compter de cette année de nouveaux moyens (délégation de courte ou longue durée) en application du décret n\textsuperscript{o} 78-216 du 2 mars 1978 relatif au régime de certaines positions des personnels enseignants titulaires de statut universitaire.

Je vous précise enfin que cette circulaire ne concerne pas les projets d'enseignement et de recherche en médecine, pharmacie, agronomie, pour lesquels les procédures sont différentes. A ce sujet, vous voudrez bien vous référer à la note ci-annexée.

\begin{flushright}
Pour le Ministre\\
et par \unsure{délégation}\\
Le \unsure{Délégué aux Relations}\\
\unsure{Universitaires Internationales}\\
\unsure{P. T.}\\
Pierre TABATONI
\end{flushright}

\margin{La présente circulaire sera publiée au Bulletin Officiel du Ministère de l'Education et du Ministère des Universités}

\hrule
\medskip
DIFFUSION, liste C, n\textsuperscript{o} 1874, NC.PB

MONTPELLIER, le 27 Mars 1979

Le Délégué aux Relations Internationales,

\unsure{F. Schué}

\uline{F. SCHUE}

\begin{flushright}
Transmis par la Cellule des Relations\\
Internationales\\
(Mme CAPELLE, poste 315) où vous pourrez\\
consulter les annexes et retirer un formulaire de présentation de projet.
\end{flushright}
